\chapter{Annexes}

\section{Dictionnaire des codes de variables EMD utilisés}
\begin{table}[H]
\centering
\caption{Principaux codes de variables EMD exploités}
\begin{tabular}{ll}
\toprule
\textbf{Code} & \textbf{Signification} \\
\midrule
I2, I4, I9 & Caractéristiques géographiques du domicile (quartier, strate, district) \\
M21 & Taille du ménage \\
M55 & Manque de transport en commun déclaré \\
M66 -- M70 & Accessibilité de base (distance, durée de marche, modes disponibles) \\
M68 & Exposition aux inondations \\
M69 & Résilience d'accès en cas de pluie \\
M73, M87 & Problèmes majeurs déclarés dans le quartier \\
M90, M91 & Accès aux écoles et daaras \\
\bottomrule
\end{tabular}
\end{table}

\section{Glossaire}
\begin{description}
\item[BI] Business Intelligence -- informatique décisionnelle.
\item[CETUD] Conseil Exécutif des Transports Urbains de Dakar.
\item[EMD] Enquête Ménages-Déplacements.
\item[ETL] Extract-Transform-Load -- processus d'extraction, transformation et chargement des données.
\item[GIMSI] Méthodologie de conception de tableaux de bord décisionnels.
\item[KPI] Key Performance Indicator -- indicateur clé de performance.
\item[RBAC] Role-Based Access Control -- contrôle d'accès basé sur les rôles.
\item[SMOTE] Synthetic Minority Over-sampling Technique -- technique de rééquilibrage de classes.
\end{description}

\section{Liste des technologies utilisées}
\begin{table}[H]
\centering
\caption{Stack technique du projet}
\begin{tabular}{ll}
\toprule
\textbf{Couche} & \textbf{Technologie} \\
\midrule
Préparation des données & Python, Pandas, Google Colab \\
ETL & Talend Open Studio \\
Entrepôt de données & PostgreSQL \\
Restitution décisionnelle & Knowage BI \\
Machine Learning & scikit-learn, joblib \\
Backend applicatif & FastAPI, Pydantic, uvicorn \\
Frontend applicatif & React.js, Recharts, Leaflet, Framer Motion \\
Gestion de projet & Scrum \\
\bottomrule
\end{tabular}
\end{table}

\section{Liste des points d'accès API}
\begin{table}[H]
\centering
\caption{Liste complète des endpoints REST exposés par l'API}
\begin{tabular}{ll}
\toprule
\textbf{Méthode} & \textbf{Endpoint} \\
\midrule
POST & /recommander \\
GET & /zones-risque \\
GET & /zones-risque/resume \\
GET & /api/anomalies/summary \\
GET & /api/anomalies/sites \\
GET & /api/anomalies/sites/\{id\} \\
GET & /api/segmentation/profils \\
POST & /predict-inaccessibility \\
GET & /api/ml/metrics \\
GET & /api/ml/features-importance \\
\bottomrule
\end{tabular}
\end{table}
